\documentclass{article}\usepackage[T1]{fontenc}\begin{document}
ForwardMixed.All.diagram proves the mixed hexagon for the actual forward comparison from the tagged family of weighted fibres to the weighted fibre of the single coproduct span. For arbitrary ordinary payload change T and base arrow a it retains SpanPull.All.comparison at Y and Z (including its componentPull-to-leafPull correction) and CoproductSpan.compare\_payload at both base objects. The target mixed face uses the whole-span cleavage built by tagging the given component cleavages and extending by the payload, not an unrelated choice; the source is the precise tagged mixed face used in \#202. The theorem includes zero and all successors, with the existing split endpoints and common universe levels. It is not the mixed diagram for the inverse comparison or the final evaluation.

Read the manuscript at 458ff36a4798e432dee5947afc402f9903a47f14, lines 929--953 (eq:telescope), then SpanPull.All.evaluationBase, SumChange.evaluationPayload and TaggedMixed.All.taggedDiagram to locate the missing square. The generic strictCell is the obstruction-solving step: map the component mixed comparison through the strict fibre tag and compose with the single span's chosen cartesian lift. Both this map and the whole-span mixed comparison factor the same image of the selected source lift by PayloadLift.square\_fac. Reindex.cancel gives equality in the strict fibre. The tag was already shown to preserve selected lifts in SpanPull.lift\_tag. The proof uses this computation on the actual lifts, not transport of a newly defined cleavage through an equivalence.

Strip the endpoint transports of Presented.square, apply strictCell and restore the transports to obtain presentedCell. Change from presentedTag to the earlier compare using compareTag; compareCell gives alignment of the mixed arrows, not yet their commutation with arbitrary comparison faces. At each length, leafAligned recovers the same component square at zero (identity span) or at a successor (positive-power cleavage). Finally comparisonTransport proves that the given forward comparison components are equality transports, using the explicit source-family correction and PullLaws.forwardApp. The payload faces here are also equality transports. IsTransport.hexagon hence reattaches these specific four arrows to the aligned mixed arrows. It does not identify arbitrary natural isomorphisms with identities.

Run the two Lean files under evidence/forward-mixed-author-fresh. DiagonalPairs uses a payload change which squares each type in a family; its image misses coordinate swap. A constant Boolean base endomorphism has no right inverse, and changes a chosen graded fibre from an empty product to an inhabited one. The strict comparison is instantiated at zero, and the forward hexagon at zero and two. TransportBoundary supplies an invertible nonidentity arrow at a single object that cannot be an equality transport. Equal middle arrows with that surrounding face fail the resulting equation. This is a countertest to weakening IsTransport, not a numerical model of the particular forward components.

The chosen inverse equivalences need not have identity units or counits. A mixed mates theorem must transfer both the base and payload faces of this forward diagram with their actual units and counits; its output must then paste with TaggedMixed.All.taggedDiagram. Base identity/composition of the inverse at fixed payload (\#201) does not give this. No new payload-unit or composition coherence, full pseudo-structure, universe widening, flattening, monad multiplication or strict algebra recovery follows here.
\end{document}
